% Module 7: Quantum Computation -- QFT, Phase Estimation, Shor's Algorithm
% Estimated slides: 35 | Duration: ~2.5 hours
% Key topics: Quantum Fourier Transform, phase estimation algorithm,
%             Shor's factoring algorithm (conceptual), RSA implications,
%             post-quantum cryptography, quantum complexity (BQP)
% Labs: QFT implementation in Qiskit; small Shor's demo

%%%%%%%%%%%%%%%%%%%%%%%%%%%%%%%%%%%%%%%%%%%%%%%%%%%%%%%%%%
\begin{frame}[fragile]\frametitle{}
\begin{center}
{\Large Module 7: Quantum Computation}\\[0.5em]
{\large QFT, Phase Estimation, and Shor's Algorithm}
\end{center}
\end{frame}

%%%%%%%%%%%%%%%%%%%%%%%%%%%%%%%%%%%%%%%%%%%%%%%%%%%%%%%%%%%
\begin{frame}[fragile]\frametitle{The Classical Discrete Fourier Transform}
\begin{itemize}
\item DFT: converts a signal from the time/position domain to the frequency domain
\item Finds the frequency components hidden in a signal
\item Classical FFT: $O(N \log N)$ for $N$ input points
\item Applications: signal processing, image compression (JPEG), solving differential equations
\item Quantum analogue: QFT computes the same transform but exponentially faster
\item QFT: $O((\log N)^2)$ gates vs classical FFT's $O(N \log N)$
\end{itemize}
% Suggested diagram: Signal in time domain -> Fourier transform -> frequency spectrum
\end{frame}

%%%%%%%%%%%%%%%%%%%%%%%%%%%%%%%%%%%%%%%%%%%%%%%%%%%%%%%%%%%
\begin{frame}[fragile]\frametitle{The Quantum Fourier Transform (QFT)}
\begin{itemize}
\item QFT: a unitary transform on $n$ qubits that computes the Discrete Fourier Transform
\item Input: a quantum state encoding $2^n$ amplitudes
\item Output: a quantum state encoding the Fourier transform of those amplitudes
\item Circuit: uses $n$ Hadamard gates + $O(n^2)$ controlled-phase gates
\item QFT is not directly useful for signal processing (can't read out all $2^n$ amplitudes)
\item But it is a crucial subroutine for phase estimation and Shor's algorithm
\end{itemize}
% Suggested diagram: QFT circuit for 3 qubits showing H gates and controlled-R gates
\end{frame}

%%%%%%%%%%%%%%%%%%%%%%%%%%%%%%%%%%%%%%%%%%%%%%%%%%%%%%%%%%%
\begin{frame}[fragile]\frametitle{QFT Circuit Structure}
\begin{itemize}
\item For $n$ qubits, the QFT circuit has depth $O(n^2)$ gates
\item Each qubit $k$ gets a Hadamard gate, then controlled-$R_k$ gates from all higher qubits
\item $R_k$ gate: phase rotation by $2\pi/2^k$
\item Final step: SWAP gates to reverse the qubit ordering
\item Approximate QFT: drop small-angle rotations for shallower circuits with small error
\end{itemize}
\begin{verbatim}
from qiskit.circuit.library import QFT
qft = QFT(num_qubits=4)
qft.draw('mpl')
\end{verbatim}
% Suggested diagram: Full 4-qubit QFT circuit; approximate vs full comparison
\end{frame}

%%%%%%%%%%%%%%%%%%%%%%%%%%%%%%%%%%%%%%%%%%%%%%%%%%%%%%%%%%%
\begin{frame}[fragile]\frametitle{Lab: Implementing QFT in Qiskit}
\begin{itemize}
\item Build QFT from scratch for 3 qubits (or use \texttt{QFT} from Qiskit library)
\item Prepare input state $|5\rangle$ = $|101\rangle$
\item Apply QFT and visualize output statevector
\item Verify by checking the classical DFT of the unit vector $e_5$
\item Extension: implement inverse QFT and confirm QFT$^{-1}$(QFT($|5\rangle$)) = $|5\rangle$
\end{itemize}
\end{frame}

%%%%%%%%%%%%%%%%%%%%%%%%%%%%%%%%%%%%%%%%%%%%%%%%%%%%%%%%%%%
\begin{frame}[fragile]\frametitle{Quantum Phase Estimation (QPE)}
\begin{itemize}
\item Goal: estimate the eigenvalue (phase) of a unitary operator $U$
\item If $U|\psi\rangle = e^{2\pi i \phi}|\psi\rangle$, QPE finds $\phi$ with high precision
\item Circuit: control register (counting qubits) + eigenstate register
\item Apply $U^{2^k}$ controlled on each counting qubit
\item Apply inverse QFT to counting register
\item Measure: gives binary approximation of $\phi$
\item Used as a subroutine in Shor's factoring and quantum simulation
\end{itemize}
% Suggested diagram: QPE circuit structure with control register, U^2^k gates, and IQFT
\end{frame}

%%%%%%%%%%%%%%%%%%%%%%%%%%%%%%%%%%%%%%%%%%%%%%%%%%%%%%%%%%%
\begin{frame}[fragile]\frametitle{Period Finding: The Key to Shor's Algorithm}
\begin{itemize}
\item Factoring an integer $N$ reduces to finding the period $r$ of $f(x) = a^x \mod N$
\item Period $r$: smallest positive integer such that $a^r \equiv 1 \pmod{N}$
\item Classical period finding: exponential in the number of bits of $N$
\item Quantum period finding using QPE: polynomial time --- this is Shor's speedup
\item Once $r$ is found, factors of $N$ can be computed classically using GCD
\item The quantum part is the period finding; everything else is classical
\end{itemize}
% Suggested diagram: Graph of f(x) = a^x mod N showing periodic pattern
\end{frame}

%%%%%%%%%%%%%%%%%%%%%%%%%%%%%%%%%%%%%%%%%%%%%%%%%%%%%%%%%%%
\begin{frame}[fragile]\frametitle{Shor's Algorithm: Overview}
\begin{itemize}
\item Input: composite integer $N$ to factor
\item Step 1 (classical): choose random $a < N$; check $\gcd(a,N)$ (if $> 1$, done)
\item Step 2 (quantum): find the period $r$ of $f(x) = a^x \mod N$ using QPE
\item Step 3 (classical): if $r$ is even and $a^{r/2} \not\equiv -1$, compute factors
\item $p = \gcd(a^{r/2} - 1, N)$ and $q = \gcd(a^{r/2} + 1, N)$
\item Repeat with different $a$ if Step 3 fails (rare)
\item Total: polynomial in $\log N$ quantum gates; exponentially faster than classical
\end{itemize}
% Suggested diagram: Shor's algorithm flowchart distinguishing classical and quantum parts
\end{frame}

%%%%%%%%%%%%%%%%%%%%%%%%%%%%%%%%%%%%%%%%%%%%%%%%%%%%%%%%%%%
\begin{frame}[fragile]\frametitle{Shor's Algorithm: A Concrete Example}
\begin{itemize}
\item Factor $N = 15$; choose $a = 7$
\item $f(x) = 7^x \mod 15$: 7, 4, 13, 1, 7, 4, 13, 1, ...
\item Period $r = 4$
\item $7^{r/2} = 7^2 = 49$; $49 \mod 15 = 4$
\item $\gcd(4-1, 15) = \gcd(3, 15) = 3$; $\gcd(4+1, 15) = \gcd(5, 15) = 5$
\item Factors: 3 and 5! ($15 = 3 \times 5$) --- correct
\item The quantum step found $r=4$ efficiently; rest is classical arithmetic
\end{itemize}
\end{frame}

%%%%%%%%%%%%%%%%%%%%%%%%%%%%%%%%%%%%%%%%%%%%%%%%%%%%%%%%%%%
\begin{frame}[fragile]\frametitle{Why Shor's Algorithm Threatens RSA}
\begin{itemize}
\item RSA relies on the difficulty of factoring large numbers (2048-4096 bits)
\item Best classical algorithm: General Number Field Sieve; super-polynomial time
\item Shor's quantum algorithm: polynomial in the number of bits
\item A quantum computer with ~4000 logical qubits could factor 2048-bit RSA
\item 4000 logical qubits $\approx$ 4 million physical qubits with current error rates
\item Current best: ~1000+ noisy physical qubits --- not there yet
\item Timeline estimate: 10-20 years to cryptographically relevant Shor's
\end{itemize}
% Suggested diagram: Resource estimate table: bits of RSA vs logical qubits vs physical qubits vs years
\end{frame}

%%%%%%%%%%%%%%%%%%%%%%%%%%%%%%%%%%%%%%%%%%%%%%%%%%%%%%%%%%%
\begin{frame}[fragile]\frametitle{Implications for Current Cryptography}
\begin{itemize}
\item Asymmetric cryptography at risk: RSA, Diffie-Hellman, Elliptic Curve (all based on factoring/discrete log)
\item Symmetric cryptography (AES-256): only quadratic speedup from Grover; double key length is sufficient
\item Hash functions (SHA-256, SHA-3): similarly, Grover's gives quadratic speedup; double output length
\item \textbf{Harvest now, decrypt later}: nation-states and adversaries collect encrypted traffic now
\item They plan to decrypt it once quantum computers are powerful enough
\item Sensitive data with 20+ year secrecy requirements should be migrated NOW
\end{itemize}
% Suggested diagram: Cryptography vulnerability matrix (algorithm type vs quantum threat level)
\end{frame}

%%%%%%%%%%%%%%%%%%%%%%%%%%%%%%%%%%%%%%%%%%%%%%%%%%%%%%%%%%%
\begin{frame}[fragile]\frametitle{Post-Quantum Cryptography: NIST Standards}
\begin{itemize}
\item NIST ran a 7-year competition (2016-2024) to standardize post-quantum algorithms
\item \textbf{CRYSTALS-Kyber} (FIPS 203): Key Encapsulation Mechanism; lattice-based
\item \textbf{CRYSTALS-Dilithium} (FIPS 204): Digital Signature; lattice-based
\item \textbf{SPHINCS+} (FIPS 205): Digital Signature; hash-based
\item \textbf{FALCON} (FIPS 206): Digital Signature; lattice-based
\item These are classical algorithms that even quantum computers cannot break (currently)
\item Migrate today: OpenSSL, TLS 1.3 implementations are being updated
\end{itemize}
% Suggested diagram: NIST PQC competition timeline and selected algorithms
\end{frame}

%%%%%%%%%%%%%%%%%%%%%%%%%%%%%%%%%%%%%%%%%%%%%%%%%%%%%%%%%%%
\begin{frame}[fragile]\frametitle{Hidden Subgroup Problem: The Unifying Framework}
\begin{itemize}
\item Shor's, Deutsch-Jozsa, Bernstein-Vazirani are all special cases of the Hidden Subgroup Problem
\item HSP: find a subgroup $H$ of a group $G$ given a function $f$ constant on cosets
\item Quantum algorithms can solve HSP for abelian groups efficiently (using QFT)
\item Non-abelian HSP (e.g., graph isomorphism) is still open --- active research
\item Unifying framework helps design new quantum algorithms for other group-theoretic problems
\end{itemize}
\end{frame}

%%%%%%%%%%%%%%%%%%%%%%%%%%%%%%%%%%%%%%%%%%%%%%%%%%%%%%%%%%%
\begin{frame}[fragile]\frametitle{Quantum Advantage Summary: What's Proven}
\begin{itemize}
\item \textbf{Exponential speedup}: Shor's (factoring), Deutsch-Jozsa, Simon's, quantum simulation
\item \textbf{Quadratic speedup}: Grover's search, amplitude estimation
\item \textbf{No speedup}: sorting, NP-complete problems (no proven exponential advantage)
\item Quantum complexity class BQP contains factoring but is not known to contain NP
\item Most quantum ``advantages'' in industry are heuristic/empirical, not proven
\item Rigorous proven advantage remains an active research frontier
\end{itemize}
% Suggested diagram: Map of problems and quantum speedup type
\end{frame}

%%%%%%%%%%%%%%%%%%%%%%%%%%%%%%%%%%%%%%%%%%%%%%%%%%%%%%%%%%%
\begin{frame}[fragile]\frametitle{Module 7 Summary}
\begin{itemize}
\item QFT: quantum Fourier transform; exponentially faster than classical FFT
\item Phase estimation: estimates eigenvalue of unitary; key subroutine
\item Shor's algorithm: factors integers in polynomial time; threatens RSA
\item Not practically feasible today: needs millions of fault-tolerant qubits
\item Post-quantum cryptography: NIST standards available now; migrate before harvest attacks
\item Proven quantum speedups: factoring (exponential), search (quadratic)
\item Next: Quantum Information --- entropy, error correction, fault tolerance
\end{itemize}
\end{frame}
